\chapter{Analyse des Besoins et Choix Technologiques}

Ce chapitre présente l'analyse complète des besoins fonctionnels et non fonctionnels qui ont guidé le développement de l'application de gestion des processus de recrutement, ainsi que les choix technologiques effectués pour répondre à ces exigences, tant pour la partie cliente que serveur.

\section{Analyse des Besoins}

L'objectif principal est de développer une application web complète pour la gestion des processus de recrutement, comprenant une interface utilisateur riche et réactive, ainsi qu'un backend robuste et sécurisé.

\subsection{Besoins Fonctionnels}

Les fonctionnalités essentielles requises pour le système sont les suivantes :

\begin{itemize}
    \item \textbf{Gestion des Utilisateurs et Authentification} : Le système doit permettre aux utilisateurs de s'enregistrer et de s'authentifier de manière sécurisée pour accéder à leurs espaces personnels. L'application doit gérer différents types d'utilisateurs avec des rôles distincts (\texttt{ADMIN}, \texttt{INTERVIEWER}, \texttt{CANDIDATE}, \texttt{HR MANAGER}), chacun disposant de permissions spécifiques. Les utilisateurs authentifiés doivent pouvoir consulter et modifier leurs informations personnelles, gérer leur mot de passe et, si nécessaire, supprimer leur compte.

    \item \textbf{Gestion du Processus de Recrutement} : Les administrateurs doivent pouvoir créer, consulter, modifier et supprimer des offres d'emploi (\texttt{Job}). Ces offres doivent être accessibles aux candidats pour consultation et candidature. Les candidats doivent pouvoir postuler aux offres d'emploi disponibles. Le système doit suivre le statut de chaque candidature (\texttt{Application}) tout au long du processus de recrutement. Les recruteurs doivent pouvoir définir et gérer leurs créneaux horaires de disponibilité (\texttt{Slot}) pour la conduite des entretiens. Le système doit permettre l'organisation d'entretiens (\texttt{Interview}) en associant un candidat, un recruteur et un créneau horaire disponible.

    \item \textbf{Interface Utilisateur et Navigation} : Après connexion, l'utilisateur doit être redirigé vers un tableau de bord personnalisé centralisant les informations et actions principales selon son rôle. L'accès à certaines pages et fonctionnalités doit être restreint selon l'état d'authentification et les permissions de l'utilisateur. L'interface doit être disponible en plusieurs langues (français et anglais initialement) avec la possibilité pour l'utilisateur de changer la langue dynamiquement.

    \item \textbf{Communication et Intégration} : L'ensemble des fonctionnalités backend doit être exposé via une API RESTful bien définie, permettant une communication claire et standardisée avec l'application cliente. L'API doit être auto-documentée pour simplifier son intégration et son utilisation par les développeurs.
\end{itemize}

\subsection{Besoins Non-Fonctionnels}

\begin{itemize}
    \item \textbf{Performance} : L'application doit être rapide et réactive, avec des temps de chargement minimaux pour garantir une expérience utilisateur fluide. Le système backend doit être capable de traiter les requêtes avec une faible latence, même sous charge. L'interface utilisateur doit offrir une navigation instantanée entre les pages grâce à l'architecture SPA.

    \item \textbf{Sécurité} : L'application doit garantir la confidentialité et l'intégrité des données utilisateur. L'accès aux ressources doit être strictement contrôlé par un mécanisme d'authentification et d'autorisation robuste. Les mots de passe doivent être stockés de manière sécurisée en utilisant des techniques de hachage et de salage. La communication entre le client et le serveur doit être sécurisée et protégée contre les attaques courantes.

    \item \textbf{Maintenabilité et Évolutivité} : L'architecture du code doit être modulaire et bien organisée pour faciliter les mises à jour futures. Le système doit être conçu pour permettre l'ajout facile de nouvelles fonctionnalités sans impact majeur sur l'existant. Le code doit suivre les bonnes pratiques et les patterns de conception reconnus.

    \item \textbf{Fiabilité et Robustesse} : Le code doit être fiable et prévisible, en minimisant les erreurs à l'exécution. L'application doit gérer les erreurs de manière gracieuse, en fournissant des retours clairs sans exposer de détails d'implémentation sensibles. Le système doit maintenir sa stabilité même en cas d'entrées invalides ou de conditions exceptionnelles.
\end{itemize}

\section{Choix Technologiques et Justifications}

Pour répondre aux besoins identifiés, une pile technologique moderne et éprouvée a été sélectionnée, combinant des technologies de pointe pour le frontend et le backend.

\subsection{Technologies Frontend}

\subsubsection{Technologies Core}

\begin{itemize}
    \item \textbf{React (v19.1.1)} : Bibliothèque JavaScript développée par Meta, choisie pour sa puissance dans la création d'interfaces utilisateur dynamiques et performantes basées sur une architecture de composants. Sa popularité garantit un large écosystème et un support communautaire solide.

    \item \textbf{TypeScript (v5.8.3)} : Sur-ensemble de JavaScript ajoutant le typage statique. Il améliore la robustesse et la maintenabilité du code en détectant les erreurs de type en amont, répondant ainsi au besoin de fiabilité.

    \item \textbf{Vite (v7.1.2)} : Outil de build nouvelle génération offrant un serveur de développement extrêmement rapide avec rechargement à chaud (HMR) et des optimisations de build pour la production, répondant au besoin de performance.

    \item \textbf{Axios (v1.11.0)} : Bibliothèque JavaScript utilisée pour effectuer des requêtes HTTP depuis le frontend React vers l'API Spring Boot. Elle permet de récupérer ou d'envoyer des données de manière asynchrone, de gérer les réponses JSON et les erreurs, répondant ainsi au besoin de communication fiable avec le serveur.
\end{itemize}

\subsubsection{Gestion de l'État et des Données}

\begin{itemize}
    \item \textbf{TanStack Query (v5.85.5)} : Bibliothèque puissante pour la gestion des données asynchrones, la mise en cache, la synchronisation et la gestion des états de serveur. Elle simplifie considérablement la gestion des requêtes API, offre des fonctionnalités avancées de cache et de synchronisation en arrière-plan, améliorant les performances et l'expérience utilisateur.

    \item \textbf{React Context API} : Solution native de React pour la gestion d'état global, utilisée notamment pour le contexte d'authentification et les toasts de notification, permettant de partager les informations utilisateur sans "prop drilling".
\end{itemize}

\subsubsection{Navigation et Routage}

\begin{itemize}
    \item \textbf{React Router DOM (v7.8.1)} : Bibliothèque de routage standard pour les applications React, choisie pour gérer la navigation côté client dans cette Single Page Application, incluant la mise en place de routes protégées nécessitant une authentification.
\end{itemize}

\subsubsection{Interface Utilisateur et Styling}

\begin{itemize}
    \item \textbf{Tailwind CSS (v4.1.12)} : Framework CSS utility-first moderne, choisi pour sa flexibilité et sa productivité. Il permet de créer rapidement des interfaces personnalisées sans écrire de CSS personnalisé, tout en maintenant une cohérence de design. L'approche utility-first accélère le développement et facilite la maintenance.

    \item \textbf{Lucide React (v0.541.0)} : Bibliothèque d'icônes SVG moderne et légère, offrant une large collection d'icônes cohérentes et personnalisables pour enrichir l'interface utilisateur.

    \item \textbf{React Icons (v5.5.0)} : Collection complète d'icônes populaires (Font Awesome, Material Design, etc.) sous forme de composants React, offrant une flexibilité maximale pour l'iconographie de l'application.

    \item \textbf{React Type Animation (v3.2.0)} : Bibliothèque spécialisée pour créer des animations de texte attractives, améliorer l'expérience utilisateur avec des effets visuels modernes et professionnels.
\end{itemize}

\subsubsection{Internationalisation}

\begin{itemize}
    \item \textbf{i18next (v25.3.6)} : Framework d'internationalisation puissant et flexible, choisi pour gérer les traductions de l'interface et permettre le changement de langue dynamique.

    \item \textbf{React i18next (v15.6.1)} : Intégration spécifique de i18next pour React, offrant des hooks et composants optimisés pour la gestion des traductions dans les composants React.
\end{itemize}

\subsubsection{Outils de Développement et Qualité du Code}

\begin{itemize}
    \item \textbf{ESLint (v9.33.0)} : Outil de linting JavaScript/TypeScript pour maintenir la qualité et la cohérence du code. Configuré avec des plugins spécifiques à React pour détecter les anti-patterns et erreurs courantes.

    \item \textbf{TypeScript ESLint (v8.39.1)} : Extension ESLint pour TypeScript, offrant des règles spécifiques au typage statique et aux bonnes pratiques TypeScript.

    \item \textbf{PostCSS (v8.5.6)} avec Autoprefixer (v10.4.21) : Outil de post-traitement CSS qui automatise l'ajout des préfixes vendeur, garantissant la compatibilité cross-browser sans effort manuel.
\end{itemize}

\subsection{Technologies Backend}

\subsubsection{Langage et Framework Principal}

\begin{itemize}
    \item \textbf{Java 21} : Dernière version LTS de Java, choisie pour ses performances optimisées, ses nouvelles fonctionnalités (Pattern Matching, Records, etc.), et sa compatibilité avec l'écosystème Spring moderne. Le support des \texttt{--enable-preview} permet d'utiliser les fonctionnalités expérimentales.

    \item \textbf{Spring Boot 3.5.4} : Framework accélérant considérablement le développement en simplifiant la configuration des applications Spring. Il intègre un serveur web embarqué, gère les dépendances efficacement et favorise les bonnes pratiques comme l'inversion de contrôle et l'injection de dépendances.
\end{itemize}

\subsubsection{Gestion des Données}

\begin{itemize}
    \item \textbf{Spring Data JPA} : Spring Data JPA simplifie l'implémentation de la couche d'accès aux données en réduisant le code répétitif. Permet l'utilisation de méthodes de requête dérivées du nom et de requêtes JPQL personnalisées.

    \item \textbf{Hibernate} : Implémentation ORM intégrée avec Spring Data JPA, permettant de mapper les objets Java aux tables de base de données de manière transparente. La configuration \texttt{hibernate.ddl-auto=update} permet la synchronisation automatique du schéma de base de données.

    \item \textbf{MySQL 8} : Base de données relationnelle robuste et performante, avec le driver \texttt{mysql-connector-j} et le dialecte \texttt{MySQL8Dialect} pour une optimisation spécifique à cette version.
\end{itemize}

\subsubsection{Sécurité et Documentation}

\begin{itemize}
    \item \textbf{Spring Security v6} : Solution de référence pour la gestion de l'authentification et de l'autorisation dans l'écosystème Spring. Configuré pour implémenter une authentification moderne avec support natif des JWT et une approche fonctionnelle.

    \item \textbf{JWT (JSON Web Tokens)} : Authentification stateless avec la bibliothèque \texttt{jjwt v0.11.5}, incluant l'API principale, l'implémentation et la sérialisation JSON avec Jackson.

    \item \textbf{Authentification par Cookies} : Système d'authentification sécurisé utilisant des cookies HTTP-only avec le filtre personnalisé \texttt{JwtCookieAuthenticationFilter}, réduisant les risques d'attaques XSS.

    \item \textbf{GlobalExceptionHandler} : Mécanisme de gestion centralisée des exceptions avec \texttt{@RestControllerAdvice}, garantissant que le client reçoit systématiquement une réponse d'erreur standardisée au format JSON.
\end{itemize}

\subsubsection{Validation et Mapping}

\begin{itemize}
    \item \textbf{Jakarta Validation 3.0.2} : Standard de validation des données avec Hibernate Validator 8.0.1.Final comme implémentation, assurant la validation automatique des DTOs et entités.

    \item \textbf{ModelMapper 3.1.1} : Bibliothèque de mapping automatique entre entités JPA et DTOs, configurée via \texttt{ModelMapperConfiguration} pour optimiser les conversions et réduire le code boilerplate.
\end{itemize}

\subsubsection{Utilitaires et Développement}

\begin{itemize}
    \item \textbf{Lombok} : Génération automatique de code (getters, setters, constructeurs, etc.) réduisant significativement le code boilerplate et améliorant la lisibilité.

    \item \textbf{Apache Commons Lang3 v3.18.0} : Bibliothèque d'utilitaires pour la manipulation de chaînes, collections et autres opérations courantes.

    \item \textbf{Spring Boot DevTools} : Outils de développement permettant le rechargement automatique de l'application lors des modifications, accélérant le cycle de développement.
\end{itemize}

\subsection{Choix des Outils}

\begin{itemize}
    \item \textbf{Git} : Système de contrôle de version distribué, essentiel pour gérer les différentes versions du code source et collaborer efficacement au sein de l'équipe de développement.

    \item \textbf{Docker} : Outil de conteneurisation permettant de créer, déployer et exécuter des applications dans des conteneurs, assurant ainsi une cohérence entre les environnements de développement, de test et de production.

    \item \textbf{Postman} : Outil utilisé pour tester les API en facilitant l'envoi de requêtes HTTP et l'analyse des réponses, utile pour le développement et la documentation de l'API.

    \item \textbf{Visual Studio Code} : Éditeur de code source léger mais puissant, avec des extensions pour le développement en Java, JavaScript et TypeScript, offrant une expérience de développement agréable et productive.

    \item \textbf{WebStorm} : IDE spécialisé de JetBrains pour le développement frontend JavaScript/TypeScript et React. Il offre des fonctionnalités avancées d'auto-complétion, de refactoring, de débogage et d'intégration avec les outils de build modernes comme Vite, optimisant le développement de l'interface utilisateur.

    \item \textbf{IntelliJ IDEA} : Environnement de développement intégré professionnel de JetBrains pour Java et Spring Boot. Il propose des outils sophistiqués pour le développement backend, incluant la gestion des dépendances Maven, l'intégration avec les bases de données, et des capacités de débogage avancées pour les applications Spring.
\end{itemize}

\section{Points d'Attention et Améliorations}

\subsection{Documentation API}

\textbf{Observation} : Le projet contient un fichier \texttt{OpenApiConfiguration.java} commenté, indiquant une intention d'intégrer Swagger/OpenAPI pour la documentation automatique de l'API, mais les dépendances correspondantes ne sont pas présentes dans le \texttt{pom.xml}.

\textbf{Recommandation} : Intégrer SpringDoc OpenAPI pour générer automatiquement la documentation interactive de l'API :

\begin{verbatim}
<dependency>
    <groupId>org.springdoc</groupId>
    <artifactId>springdoc-openapi-starter-webmvc-ui</artifactId>
    <version>2.3.0</version>
</dependency>
\end{verbatim}

\section{Justification de l'Architecture Globale}

L'ensemble des choix technologiques effectués forme une architecture cohérente et moderne qui répond parfaitement aux besoins identifiés :

\begin{itemize}
    \item \textbf{Séparation des préoccupations} : L'architecture client-serveur avec API RESTful assure une séparation claire entre la logique de présentation et la logique métier.

    \item \textbf{Scalabilité} : La nature découplée de l'architecture permet de faire évoluer indépendamment le frontend et le backend selon les besoins.

    \item \textbf{Expérience développeur} : Les outils modernes comme Vite et Spring Boot offrent une excellente expérience de développement avec des temps de compilation rapides et une configuration simplifiée.

    \item \textbf{Sécurité renforcée} : L'utilisation de JWT avec Spring Security côté serveur et la gestion d'état sécurisée côté client garantissent un niveau de sécurité élevé.

    \item \textbf{Maintenabilité à long terme} : Le typage fort (TypeScript et Java) et l'architecture modulaire facilitent la maintenance et l'évolution de l'application.

    \item \textbf{Performance optimisée} : L'utilisation de TanStack Query pour la gestion des données asynchrones et Tailwind CSS pour le styling améliore considérablement les performances et l'expérience utilisateur.

    \item \textbf{Interface utilisateur moderne} : La combinaison de React 19, Tailwind CSS et des bibliothèques d'animation permet de créer une interface utilisateur moderne, responsive et attrayante.
\end{itemize}

Cette combinaison technologique représente un équilibre optimal entre performance, sécurité, maintenabilité et expérience utilisateur, assurant ainsi le succès du projet à court et long terme.
